% !TEX root = ../../../main.tex
% 所属：第二部分：多变量分析学 / 第2章：多元函数微分学
% 节：2.3 中值定理与 Taylor 公式
% 正文按原文件顺序保存；由 ../chapter2.tex 调用。
\section{中值定理与 Taylor 公式}

\subsection*{中值定理}

\begin{definition}
	设 $D\subset\R^n$.若对任意两点 $\vect x,\vect y\in D$ 及任意
	$\lambda\in[0,1]$，均有
	\[
	\lambda\vect x+(1-\lambda)\vect y\in D,
	\]
	则称 $D$ 为\textbf{凸集}.换言之，$D$ 中任意两点之间的线段都包含在 $D$ 中.
\end{definition}

\begin{theorem}[Lagrange 中值定理]
	设 $D\subset\R^2$ 为凸开集，函数 $f:D\to\R$ 在 $D$ 上可微.
	对任意 $(x_1,y_1),(x_2,y_2)\in D$，存在 $\theta\in(0,1)$，使得
	\[
	\begin{aligned}
	&f(x_1,y_1)-f(x_2,y_2)\\
	&\quad=
	f_x\bigl(x_1+\theta(x_2-x_1),y_1+\theta(y_2-y_1)\bigr)(x_1-x_2)\\
	&\qquad+
	f_y\bigl(x_1+\theta(x_2-x_1),y_1+\theta(y_2-y_1)\bigr)(y_1-y_2).
	\end{aligned}
	\]
\end{theorem}

\begin{proof}
	令
	\[
	\varphi(t)
	=
	f\bigl(x_1+t(x_2-x_1),y_1+t(y_2-y_1)\bigr),
	\qquad t\in[0,1].
	\]
	由于 $D$ 为凸集，$\varphi$ 在 $[0,1]$ 上有定义.由 $f$ 可微，
	$\varphi$ 在 $[0,1]$ 上连续并在 $(0,1)$ 上可导，且
	\[
	\begin{aligned}
	\varphi'(t)
	={}&
	f_x\bigl(x_1+t(x_2-x_1),y_1+t(y_2-y_1)\bigr)(x_2-x_1)\\
	&+
	f_y\bigl(x_1+t(x_2-x_1),y_1+t(y_2-y_1)\bigr)(y_2-y_1).
	\end{aligned}
	\]
	对 $\varphi$ 应用一元函数的 Lagrange 中值定理，存在 $t_0\in(0,1)$，使得
	\[
	\varphi(1)-\varphi(0)=\varphi'(t_0).
	\]
	将上式整理，并令 $\theta=t_0$，即得所证公式.
\end{proof}

\begin{proposition}
	设 $D\subset\R^2$ 为区域，函数 $f\in C^1(D)$.若在 $D$ 上
	\[
	\frac{\partial f}{\partial x}=0,
	\qquad
	\frac{\partial f}{\partial y}=0,
	\]
	即 $\operatorname{grad}f=\vect 0$，则 $f$ 在 $D$ 上为常数.
\end{proposition}

\begin{proof}
	若 $D$ 为凸区域，则对任意 $P,Q\in D$，由 Lagrange 中值定理立即得到
	\[
	f(P)-f(Q)=0,
	\]
	故 $f(P)=f(Q)$.

	一般地，区域 $D$ 是开连通集，因而是折线连通的.对任意 $P,Q\in D$，
	存在完全包含在 $D$ 中的折线
	\[
	P=P_0,P_1,\ldots,P_m=Q.
	\]
	每条线段 $[P_{i-1},P_i]$ 都包含在 $D$ 中.
	在各条线段上依次应用 Lagrange 中值定理，有
	\[
	f(P_{i-1})=f(P_i),
	\qquad i=1,2,\ldots,m.
	\]
	于是
	\[
	f(P)=f(P_1)=\cdots=f(P_m)=f(Q).
	\]
	由于 $P,Q$ 任意，故 $f$ 在 $D$ 上为常数.
\end{proof}

\begin{theorem}[多元函数的中值公式]
	设 $n$ 元函数 $f(x_1,x_2,\ldots,x_n)$ 在凸区域
	$D\subset\R^n$ 上可微. 对于 $D$ 内任意两点
	\[
	(x_1^0,x_2^0,\ldots,x_n^0)
	\quad\text{与}\quad
	(x_1^0+\Delta x_1,x_2^0+\Delta x_2,\ldots,x_n^0+\Delta x_n),
	\]
	至少存在一个 $\theta\in(0,1)$, 使得
	\[
	\begin{aligned}
	&f(x_1^0+\Delta x_1,x_2^0+\Delta x_2,\ldots,x_n^0+\Delta x_n)
	-f(x_1^0,x_2^0,\ldots,x_n^0)\\
	&\qquad=
	\sum_{i=1}^n
	f_{x_i}(x_1^0+\theta\Delta x_1,x_2^0+\theta\Delta x_2,
	\ldots,x_n^0+\theta\Delta x_n)\Delta x_i.
	\end{aligned}
	\]
	
	使用向量与点积，可以将上述的定理表述为如下形式
	设 $n$ 元函数 $f(\vect x)$ 在凸区域
	$D\subset\R^n$ 上可微. 对于 $D$ 内任意两点
	\[
	\vect x
	\quad\text{与}\quad
	\vect x +\vect{\Delta x},
	\]
	
	至少存在一个 $\theta\in(0,1)$, 使得
	\[f(\vect x +\vect{\Delta x})
		-f(\vect x )=\nabla f(\vect x +\theta \vect{\Delta x}) \cdot \vect {\Delta x}\]
\end{theorem}

\subsection*{二元函数的 Taylor 公式}
\begin{proposition}
设 $z=f(x,y)$ 在 $(x_0,y_0)$ 的某邻域内具有直到 $k+1$ 阶的连续偏导数,
且 $(x_0+\Delta x,y_0+\Delta y)$ 仍在该邻域内. 则
\[
\begin{aligned}
	f(x_0+\Delta x,y_0+\Delta y)
	={}&f(x_0,y_0)
	+\left(
	\Delta x\frac{\partial}{\partial x}
	+\Delta y\frac{\partial}{\partial y}
	\right)f(x_0,y_0)\\
	&+\cdots+
	\frac{1}{k!}
	\left(
	\Delta x\frac{\partial}{\partial x}
	+\Delta y\frac{\partial}{\partial y}
	\right)^k f(x_0,y_0)
	+R_k,
\end{aligned}
\]
其中
\[
R_k=
\frac{1}{(k+1)!}
\left(
\Delta x\frac{\partial}{\partial x}
+\Delta y\frac{\partial}{\partial y}
\right)^{k+1}
f(x_0+\theta\Delta x,y_0+\theta\Delta y),
\qquad
\theta\in(0,1),
\]
称为 Lagrange 余项.   
\end{proposition}


\begin{proof}
	令
	\[
	\varphi(t)=f(x_0+t\Delta x,y_0+t\Delta y).
	\]
	于是
	\[
	\varphi(0)=f(x_0,y_0),
	\qquad
	\varphi(1)=f(x_0+\Delta x,y_0+\Delta y).
	\]
	由于对 $t\in[0,1]$, $\varphi(t)$ 有定义且具有 $k+1$ 阶导数,
	对 $\varphi(t)$ 在 $t=0$ 处应用一元函数的 Taylor 公式,
	\[
	\varphi(1)
	=
	\varphi(0)+\varphi'(0)
	+\frac{1}{2!}\varphi''(0)
	+\cdots+
	\frac{1}{k!}\varphi^{(k)}(0)
	+\widetilde R_k,
	\]
	其中
	\[
	\widetilde R_k
	=
	\frac{1}{(k+1)!}\varphi^{(k+1)}(\theta),
	\qquad
	\theta\in(0,1).
	\]
	又
	\[
	\varphi^{(j)}(t)
	=
	\left(
	\Delta x\frac{\partial}{\partial x}
	+\Delta y\frac{\partial}{\partial y}
	\right)^j
	f(x_0+t\Delta x,y_0+t\Delta y),
	\]
	代入即得所述 Taylor 公式与 Lagrange 余项.
\end{proof}

若
\[
r=\sqrt{\Delta x^2+\Delta y^2}\to0,
\]
则还可将余项写成
\[
R_k=o(r^k)
=o\left(\left(\sqrt{\Delta x^2+\Delta y^2}\right)^k\right),
\]
称为 Peano 余项.

\begin{theorem}[$n$ 元函数的 Taylor 公式]
	设 $n$ 元函数 $f(x_1,x_2,\ldots,x_n)$ 在点
	$(x_1^0,x_2^0,\ldots,x_n^0)$ 附近具有 $k+1$ 阶的连续偏导数,
	则在该点附近成立
	\[
	\begin{aligned}
	&f(x_1^0+\Delta x_1,x_2^0+\Delta x_2,\ldots,x_n^0+\Delta x_n)\\
	&\quad=
	f(x_1^0,x_2^0,\ldots,x_n^0)
	+\left(
	\sum_{i=1}^n\Delta x_i\frac{\partial}{\partial x_i}
	\right)f(x_1^0,x_2^0,\ldots,x_n^0)\\
	&\qquad+
	\frac{1}{2!}
	\left(
	\sum_{i=1}^n\Delta x_i\frac{\partial}{\partial x_i}
	\right)^2 f(x_1^0,x_2^0,\ldots,x_n^0)
	+\cdots\\
	&\qquad+
	\frac{1}{k!}
	\left(
	\sum_{i=1}^n\Delta x_i\frac{\partial}{\partial x_i}
	\right)^k f(x_1^0,x_2^0,\ldots,x_n^0)
	+R_k,
	\end{aligned}
	\]
	其中
	\[
	\begin{aligned}
	R_k
	={}&
	\frac{1}{(k+1)!}
	\left(
	\sum_{i=1}^n\Delta x_i\frac{\partial}{\partial x_i}
	\right)^{k+1}\\
	&\quad\cdot
	f(x_1^0+\theta\Delta x_1,x_2^0+\theta\Delta x_2,
	\ldots,x_n^0+\theta\Delta x_n),
	\qquad 0<\theta<1,
	\end{aligned}
	\]
	为 Lagrange 余项.
\end{theorem}

